\documentclass[12pt, a4paper]{report}

% Packages
\usepackage[utf8]{inputenc}
\usepackage[T1]{fontenc}
\usepackage[english]{babel}
\usepackage{amsmath, amssymb, amsthm, amsfonts}
\usepackage{geometry}
\usepackage{hyperref}
\usepackage{bookmark}
\usepackage{enumitem}
\usepackage{fancyhdr}
\usepackage{titlesec}
\usepackage{mathrsfs}
\usepackage{xcolor}
\usepackage{tikz}
\usepackage{pgfplots}
\pgfplotsset{compat=1.18}
\usepackage{eurosym}

% Geometry setup
\geometry{
    top=2.5cm,
    bottom=2.5cm,
    left=2.5cm,
    right=2.5cm
}
\setlength{\headheight}{30pt}

% Hyperref setup
\hypersetup{
    colorlinks=true,
    linkcolor=blue,
    filecolor=magenta,      
    urlcolor=cyan,
    pdftitle={Stochastic Finance and Risk Modelling - Chapter 8},
    pdfauthor={Laurence Carassus and Gaoyue Guo},
}

% Theorem environments
\newtheorem{theorem}{Theorem}[chapter]
\newtheorem{proposition}[theorem]{Proposition}
\newtheorem{lemma}[theorem]{Lemma}
\newtheorem{corollary}[theorem]{Corollary}

\theoremstyle{definition}
\newtheorem{definition}{Definition}[chapter]
\newtheorem{example}{Example}[chapter]
\newtheorem{remark}[theorem]{Remark}

% Layout/Style
\pagestyle{fancy}
\fancyhf{}
\rhead{\leftmark}
\lhead{}
\cfoot{\thepage}

\title{\textbf{Stochastic Finance and Risk Modelling} \\ \large Chapter 8: Risk Measures}
\author{Based on slides by Andrea Molent \\ Laurence Carassus and Gaoyue Guo \\ \textit{Expanded Study Resource}}
\date{\today}

\begin{document}

\maketitle

\tableofcontents

\chapter*{Introduction to the Study Guide}
\addcontentsline{toc}{chapter}{Introduction to the Study Guide}
This document serves as a comprehensive study guide for the course "Stochastic Finance and Risk Modelling," functioning as a complete substitute for the lecture slides related to Chapter 8: "Risk Measures."

Proper measurement and management of risk are fundamental to the operation of financial markets, particularly highlighted by the frequent market crises that demonstrate the imperfect nature of traditional risk metrics. The objective of this chapter is to present a formal background on how to quantify and mathematically represent risk.

The discussion begins with an overview of the qualitative kinds of risks that exist, specifically zooming in on \textit{Market Risk}. We will recall necessary probabilistic concepts such as distribution functions and quantiles. Then, we establish the mathematical axioms that govern optimal and \textit{coherent} definitions of a risk measure, such as translation invariance, sub-additivity, homogeneity, and monotonicity. Following this framework, we analyze traditional risk measures such as Variance, the widely adopted \textit{Value at Risk (VaR)}, and its superior coherent alternative, the \textit{Expected Shortfall (ES)}. Finally, we will touch upon the optimal management strategy for limiting risk with regards to derivative products in varying frameworks.

Detailed explanations, rigorous mathematical derivations, and intuition are provided throughout to bridge conceptual gaps and ensure a deep understanding of the material without the necessity of external resources.

\chapter{Introduction to Financial Risks}

The landscape of modern finance has been frequently punctuated by severe losses suffered by financial institutions and industrial companies alike. These events, notably the Global Financial Crisis of 2007-2009, have emphatically shown the necessity for all participants in financial markets—ranging from intermediaries to regulatory and supervisory authorities—to deeply understand the risks intertwined with their financial activities. 

Mere awareness of risk is insufficient; it needs to be strictly \textbf{measured, monitored, and managed}. The rigor and quality of internal control within any institution depend substantially on a robust and mathematically sound understanding of these risks. The crisis highlighted that traditional risk measures are often poorly suited for times of extreme distress, necessitating more robust metrics to assess potential disaster scenarios.

\section{Taxonomy of Risks in Financial Institutions}

The activities of financial institutions are exposed to several classes of risk. According to standard definitions (e.g., Jorion, 2007), these can be primarily divided into four categories:

\begin{enumerate}
    \item \textbf{Market Risk:} This is the focal point of the current chapter. Market risk arises from unfavorable variations in macroeconomic variables that directly dictate asset pricing. This includes shocks to equity prices, fluctuating interest rates, foreign exchange rate movements, varying market indices, shifts in volatilities, and changing correlation structures among varying assets. 
    \begin{itemize}
        \item \textit{Example:} A sudden crash in equity prices directly depletes the mark-to-market value of an investment portfolio containing widespread stocks.
    \end{itemize}
    
    \item \textbf{Liquidity Risk:} The risk stemming from an institution's inability to fulfill its short-term financial obligations. This normally is not caused by fundamental insolvency, but rather poor synchronization of cash inflows and outflows (funding liquidity risk), or the inability to execute a transaction at current market prices due to lack of market depth (market liquidity risk).
    \begin{itemize}
        \item \textit{Example:} A classic "bank run," where a myriad of customers simultaneously attempt to withdraw deposits, and the bank does not have enough immediate cash on hand to service them, despite owning lucrative long-term assets.
    \end{itemize}

    \item \textbf{Credit Risk:} This involves the possibility of loss due to a counterparty's failure to meet their contractual obligations, specifically regarding the payment of interest or the restitution of principal.
    \begin{itemize}
        \item \textit{Example:} A corporate bond issuer defaulting on its obligations, or an insolvent bank failing to return funds deposited by a client.
    \end{itemize}

    \item \textbf{Operational Risk:} The hazard of loss resulting from inadequate or failed internal processes, human error, system failures, or external events (fraud, cyber attacks, legal risks).
    \begin{itemize}
        \item \textit{Example:} A "fat-finger" error where a trader accidentally inputs an extra zero, purchasing ten times the nominal amount of securities intended, at incorrectly computed limit prices.
    \end{itemize}
\end{enumerate}

\section{Focus on Market Risk}

We narrow our focus down to Market Risk, as its direct derivation from asset prices makes it particularly receptive to stochastic and probabilistic modeling. Market risk inherently differs depending on the precise nature of the underlying asset targeted:

\begin{itemize}
    \item \textbf{Currency Risk (Foreign Exchange Risk):} This arises from exposure to fluctuations in the exchange rates between currencies. It heavily affects multinational corporations that hold raw materials or revenues in foreign denominations.
    \item \textbf{Interest Rate Risk:} Variations in the interest rate structure drastically affect any agent engaged in lending, borrowing, or holding fixed-income instruments. Bonds specifically have prices that move inversely to the interest rate.
    \item \textbf{Equity Risk:} Directly relates to holding stocks. The price of an individual share behaves stochastically, influenced by the specifics of the issuing firm and encompassing macroeconomic trends.
    \item \textbf{Other Residual Risks:} This captures secondary exposure metrics. A prime example is \textbf{Volatility Risk}, which drastically affects traders of non-linear derivatives (such as options), as option prices scale with the volatility assumption (Vega exposure). Another form is \textbf{Basis Risk}, which frequently transpires in hedging mechanisms. It describes the risk incurred when the hedging instrument and the hedged asset do not behave identically (e.g., trying to hedge corporate bond exposure by shorting risk-free sovereign bonds, leaving the trader exposed to the corporate credit spread).
\end{itemize}

Furthermore, market risk can be monitored under two primary paradigms:
\begin{itemize}
    \item \textbf{Absolute Risk:} Evaluating potential losses purely in fixed denomination terms (e.g., estimating a potential portfolio drop of \$5,000,000 over ten days). This looks into the actual probability distribution of total underlying returns.
    \item \textbf{Relative Risk:} Evaluating potential losses relative to an underlying benchmark index. In this framing, the absolute return matters less than the \textit{tracking error}. An institutional investor might suffer absolute losses, but if these losses are lesser than those of the benchmark, the relative risk profile handles it as an outperformance.
\end{itemize}

Institutions manage their market risk profiles using fixed limits on notionals, gross and net exposure caps, Value at Risk (VaR) ceilings, and robust supervision facilitated by independent risk management units.

\section{The Practical Necessity of Risk Measurement}

Carefully estimating the risk landscape surrounding the use of modern financial instruments is absolutely vital for several pragmatic purposes:

\begin{enumerate}
    \item \textbf{Optimal Allocation of Regulatory Capital:} Equity is a heavily constrained and expensive resource for banks. Financial institutions must reserve sufficient capital (often termed Tier 1 or Tier 2 capital) proportionate to the riskiness of their holdings as a bankruptcy buffer. Precise risk measures prevent the unnecessary hoarding of idle capital.
    \item \textbf{Risk Monitoring and Hedging:} Having a quantified mathematical measure enables the construction of limits (like "traffic light" systems—green for safe, yellow for cautionary, red for immediate liquidation limits) for trading desks, specifically targeting highly leveraged, illiquid, or complex derivative assets.
    \item \textbf{Compliance:} Conformity with strict guidelines issued by national banking authorities, financial supervisory committees, and international directives cannot be maintained without standardized mechanisms for computing theoretical downside variations.
\end{enumerate}

\section{The Regulatory Environment}

Financial regulations function concurrently across multiple layers:

\begin{itemize}
    \item \textbf{National Level:} Central banks, financial stability boards, and domestic clearinghouses impose stringent checks to assure systemic stability concerning the activities within their specific jurisdictions.
    \item \textbf{International Level (Basel Committee):} Beginning with its "Capital Adequacy Directive" of 1995, the \textbf{Basel Committee on Banking Supervision} advocated the widespread employment of formal methods for quantifying market risk. Their approach fundamentally endorsed the employment of the \textbf{Value at Risk (VaR)} metric through banks' sophisticated internal models.
\end{itemize}

Subsequent iterations, notably the \textbf{Basel II} frameworks (2004) and more recently \textbf{Basel III/IV}, have progressively refined these directives. The recent revisions heavily incorporated lessons extracted from the 2007-2009 collapse, aiming to curb excessive leverage and require capitalization against not just standardized VaR, but also Stressed VaR and broader, cohesive Expected Shortfall metrics that catch "tail events."

\section{Primer on Value at Risk (VaR) and Associated Critiques}

Driven to enormous popularity by J.P. Morgan in the early 1990s through their "RiskMetrics" system, the \textbf{Value at Risk (VaR)} endeavors to condense the holistic risk affecting a massively complex, multi-asset portfolio into a simple, single, explainable cash value understandable by standard executives.

At a given statistical \textit{confidence interval} (such as $95\%$ or $99\%$), VaR attempts to quantify the maximum potential amount of loss an institution could face over a prescribed short period (for instance, $1$ day or $10$ days) assuming "normal" market conditions. The psychological ancestry of VaR links back to the "probability of ruin" theory utilized heavily by actuarial scientists and insurance firms.

While theoretically simple at heart—interpretable precisely as a specific \textit{quantile} of an estimated loss distribution—its calculation procedure rapidly balloons in complexity when assessing highly non-linear products (like path-dependent exotic options) or dealing with unpredictable empirical correlations between thousands of assets in a global portfolio.

Despite its industry dominance, the VaR possesses monumental theoretical and practical flaws:
\begin{itemize}
    \item It disregards entirely the magnitude of losses \textit{exceeding} the specific quantile target. Thus, a portfolio holding an asset that causes a \$1 million loss with 1\% probability is viewed inversely identical to one causing a \$100 billion loss with 1\% probability via the lens of a $99\%$ VaR.
    \item It is often mathematically blind to diversification benefits, violating basic theoretical axioms of risk evaluation.
\end{itemize}

These profound deficiencies catalyzed the widespread development and endorsement of superior indicators—principally the \textbf{Expected Shortfall (ES)}—which we will meticulously explore in coming sections.

\section{Key Investor Information Document (KIID) and Retail Communication}

Risk measures are not strictly tools nested within the complex quantitative desks of investment banks; they are a direct tool for retail client communication as mandated by European regulatory frameworks (such as UCITS and PRIIPs).

A typical retail fund is obligated to issue a \textit{Key Investor Information Document (KIID)}. The KIID exhibits a visually prominent \textbf{Synthetic Risk and Reward Indicator (SRRI)}, categorizing the fund upon a scale from 1 (lowest risk/reward) to 7 (highest risk/reward).

\textit{Note on Interpretation:} The KIID graphic utilizes historically captured annualized volatility over trailing periods (commonly 5 years). A score of 1 does not infer "guaranteed safety," but rather extremely compressed volatility bands.

The KIID will also append vital qualitative warnings reflecting risks that synthetic, historical quantiles inherently fail to project:
\begin{itemize}
    \item \textbf{Credit Event Risk:} Sudden multi-notch downgrade of counterparties.
    \item \textbf{Liquidity Freezes:} An abrupt evaporation of quoting volume, forcing substantial pricing hits if liquidation is forcefully necessary.
    \item \textbf{Counterparty Failure:} In specifically structured swaps or total return structures.
    \item \textbf{Emerging Markets Profile:} Emphasizing unique external constraints such as systemic political instability or unrefined legal recourse structures, which mathematically escape simple pricing variance metrics computed over benign economic epochs.
\end{itemize}

\chapter{Probability Reminders: Distribution Functions and Quantiles}

Before defining formal risk measures, it is necessary to rigorously establish the mathematical foundations regarding the distribution of random variables. These concepts lie at the heart of evaluating potential losses and mathematically describing exposure.

\section{Cumulative Distribution Functions (CDF)}

\begin{definition}[Cumulative Distribution Function]
Let $X$ be a real-valued random variable defined on a probability space $(\Omega, \mathcal{F}, \mathbb{P})$. We term the \textbf{Cumulative Distribution Function (CDF)} of $X$ as the function $F_X : \mathbb{R} \to [0, 1]$ defined by:
$$
\forall x \in \mathbb{R}, \quad F_X(x) = \mathbb{P}(X \leq x).
$$
\end{definition}

The cumulative distribution function acts as a complete probabilistic descriptor for a random variable. A defining characteristic of the CDF is its adherence to the following proposition:

\begin{proposition}[Properties of the CDF]
The cumulative distribution function $F_X$ of a random variable $X$ possesses the following properties:
\begin{enumerate}
    \item \textbf{Weakly increasing:} If $x_1 \leq x_2$, then $F_X(x_1) \leq F_X(x_2)$.
    \item \textbf{Càdlàg (Continu à droite, limite à gauche):} Bounded and right-continuous. Specifically, for every point $x_0 \in \mathbb{R}$:
    $$
    \lim_{x \to x_0^+} F_X(x) = F_X(x_0).
    $$
    It has left limits at every point, meaning $\lim_{x \to x_0^-} F_X(x)$ exists, and represents $\mathbb{P}(X < x_0)$.
    \item \textbf{Asymptotic limits:} It behaves predictably at infinities:
    $$
    \lim_{x \to -\infty} F_X(x) = 0, \quad \lim_{x \to +\infty} F_X(x) = 1.
    $$
\end{enumerate}
\end{proposition}

\subsection*{Example: The Standard Normal Distribution}

A crucial distribution in financial modeling is the standard normal distribution $\mathcal{N}(0, 1)$, whose density function is typically denoted as $\phi(x)$ and CDF as $\Phi(x)$. While $\Phi(x)$ lacks a closed-form algebraic formula, its properties are widely documented. Notably:
\begin{itemize}
    \item When addressing tail risks, critical reference points are frequently encountered. The $95^{\text{th}}$ percentile yields $\Phi(1.64) \approx 0.95$.
    \item Moving to the deep tail, the $99^{\text{th}}$ percentile gives $\Phi(2.33) \approx 0.99$.
\end{itemize}
These specific critical values directly translate into standard Value at Risk calculations.

\subsection{Characterizations and Limit Theorems}

The CDF has expansive theoretical applications inside probability theory:
\begin{enumerate}
    \item \textbf{Law Characterization:} The CDF characterizes the law of the random variable. If two random variables $X$ and $Y$ share the exact same CDF, they possess the exact same distribution. Thus, for any bounded measurable function $f: \mathbb{R} \to \mathbb{R}$:
    $$
    \mathbb{E}[f(X)] = \mathbb{E}[f(Y)].
    $$
    \item \textbf{Weak Convergence:} A sequence of random variables $(X_n)_{n \geq 1}$ is said to converge in law (or weak convergence) to a limit $X$ if and only if, for every continuity point $x \in \mathbb{R}$ of the specific function $F_X$, the point-wise subsequent limit holds:
    $$
    \lim_{n \to +\infty} F_{X_n}(x) = F_X(x).
    $$
    \item \textbf{Empirical Approximation (Glivenko-Cantelli Theorem):} Let an $i.i.d.$ sequence of random variables $(X_n)_{n \geq 1}$ be given. We can construct an empirical cumulative distribution function $\hat{F}_n(x) = \frac{1}{n} \sum_{i=1}^n \mathbf{1}_{\{X_i \leq x\}}$. The fundamental Glivenko-Cantelli Theorem guarantees that this empirical distribution converges almost surely to the true distribution uniformly over $\mathbb{R}$:
    $$
    \sup_{x \in \mathbb{R}} |\hat{F}_n(x) - F_X(x)| \xrightarrow{a.s.} 0.
    $$
    This is deeply critical in quantitative finance, as real-world data distributions (like historical asset returns) are usually modeled by directly leaning on empirical distribution functions.
\end{enumerate}

\section{Quantiles and the Generalized Inverse}

Strict inversion of a CDF is frequently impossible because a CDF may not be strictly increasing (it might have flat "plateaus" where probabilities are zero) and it may not be continuous (exhibiting "jumps" where a discrete point probability mass resides). To bypass this complication, we resort to the concept of the \textit{generalized inverse}.

\begin{definition}[Quantiles and Generalized Inverse]
Let $X$ be a random variable alongside its CDF $F_X$. For any probability level $p \in (0, 1]$, we define the \textbf{quantile of order $p$} (or $p$-th quantile) as the threshold value:
$$
x_p = \inf \{x \in \mathbb{R} : F_X(x) \geq p\}.
$$
By mathematical convention, $\inf\{\emptyset\} = +\infty$.

We define the function:
\begin{align*}
    F_X^{-1} : (0, 1) &\to \mathbb{R} \\
    p &\mapsto x_p
\end{align*}
This function is right-continuous and is formally termed the \textbf{generalized inverse} of $F_X$.
\end{definition}

In practical terminology based on the earlier standard normal example, we deduce that $\Phi^{-1}(0.95) \approx 1.64$ equating to the quantile $x_{0.95} \approx 1.64$, and $\Phi^{-1}(0.99) \approx 2.33$, meaning $x_{0.99} \approx 2.33$.

\subsection{Properties of the Generalized Inverse}

The interplay between $F_X$ and $F_X^{-1}$ is governed by specific structured boundaries.

\begin{proposition}[Inequality Relations]
Let $F_X$ be a valid CDF, and $F_X^{-1}$ be its generalized inverse. For any generic probability $p \in (0, 1)$, the following lower-bound statement holds:
$$
F_X\left(F_X^{-1}(p)\right) \geq p.
$$
Moreover, the generalized inverse satisfies the fundamental equivalence:
$$
F_X(x) \geq p \iff x \geq F_X^{-1}(p).
$$
\end{proposition}
\begin{remark}
The inequality $F_X(F_X^{-1}(p)) \geq p$ arises precisely because of potential "jumps" (discontinuities) in the CDF $F_X$. If $F_X$ is completely continuous over $\mathbb{R}$, this inequality compresses into a clean equality:
$$
\forall p \in (0, 1), \quad F_X(F_X^{-1}(p)) = p.
$$
\end{remark}

\subsection*{Stochastic Simulation and Generating Variables}
The generalized inverse holds massive importance for \textbf{Monte Carlo simulations} frequently utilized by financial engineers to price complex derivatives.

\begin{itemize}
    \item \textbf{The Inverse Transform Theorem:} Suppose $U$ is a random variable following a uniform distribution over $[0, 1]$. Then, the custom-generated random variable structurally defined as $Z = F_X^{-1}(U)$ holds identically the exact same law as the target random variable $X$. This means an engine that dynamically generates uniform random variables can simulate theoretically \textit{any} probability distribution purely by feeding the uniform variables into that distribution's generalized inverse.
    \item Conversely, if $F_X$ is strictly continuous, plotting the transformation $U' = F_X(X)$ yields a random variable explicitly following a uniform distribution within $[0, 1]$.
\end{itemize}

Ultimately, utilizing the generalized inverse is essential precisely because it correctly formalizes quantiles even upon probability spaces filled with complex, mixed continuous-discrete dynamics (such as the default probability of discrete bond issuances intertwined with continuous systemic stock indices).

\chapter{Axiomatic Definitions and Properties of Risk Measures}

To navigate risk analytically, we must transition from vague, abstract notions of "danger" toward an exact, measurable numeric representation. The mathematical framework developed extensively over the previous two decades provides an axiomatic set of conditions that an objectively "good" risk measure should satisfy. 

\section{Framework and Notation}

We focus primarily on addressing the risk affecting the mark-to-market value of a specified long position. Let us formulate the base variables upon a complete measurable space $(\Omega, \mathcal{F})$. Then we define:
\begin{itemize}
    \item $V_0$: the currently known deterministic value (at time $t=0$) of the portfolio.
    \item $V_h$: the stochastic random variable outlining the terminal value of the portfolio evaluated precisely at a continuous horizon $h$ (for example, $h = 10$ trading days).
    \item $L := V_0 - V_h$: the theoretical \textbf{loss} between $t=0$ and terminal horizon $t=h$. Thus, $L$ acts as a random variable where a strictly positive realization represents a pure monetary loss, whilst a negative realization mathematically signifies a monetary gain.
\end{itemize}

The fundamental danger inherent in the portfolio is dictated totally by the probabilistic spectrum of $L$.
By denoting its Cumulative Distribution Function as $F_L(x) = \mathbb{P}(L \leq x)$, we attempt to construct an indicator function, called a \textit{risk measure}, commonly dubbed $\rho(L)$, that condenses the portfolio's entire riskiness profile into a singular, unambiguous scalar number. 

Trivially, the "riskier" an entire position behaves structurally—involving higher likelihoods of catastrophic $L > 0$ realizations—the numerically higher $\rho(L)$ must be. We can contextualize the output of any given generic $\rho(L)$ functionally via absolute capital requirements:
\begin{enumerate}
    \item $\rho(L) > 0$: The portfolio is overly risky. $\rho(L)$ fundamentally signifies the exact amount of buffer cash capital the participating agent \textbf{must inject} directly alongside the risky position to downgrade its risk level such that it behaves as a perfectly acceptable position per internal protocols.
    \item $\rho(L) = 0$: The portfolio behaves at the precise equilibrium of accepted boundaries. It flawlessly satisfies immediate capital risk requisites without excess cushion.
    \item $\rho(L) < 0$: The portfolio behaves exceptionally safely. The output $-\rho(L)$ structurally represents the surplus monetary quantity that can be cleanly \textbf{extracted} away from the position without violating capital requirement stipulations, allowing that cash to be invested into higher-yielding ventures.
\end{enumerate}

\section{Preliminary Examples of Proposed Risk Measures}

We constrain our analysis exclusively to bounded financial scenarios. We define $L^\infty$ as the constrained set of theoretical positions $L$ strictly conforming to the bound:
$$
\|L\|_\infty = \sup_{\omega \in \Omega} |L(\omega)| < +\infty.
$$

Let any general function $\rho: A \subseteq L^\infty \to \mathbb{R}$ dictate our "risk". Across classical financial literature, numerous simplistic attempts have populated risk departments:

\begin{enumerate}
    \item \textbf{Maximum Loss Method:} $\rho_{\max}(L) = \sup_{\omega \in \Omega} L(\omega)$. This evaluates purely the worst-case imaginable scenario. Unfortunately, this forms a horribly terrible implementation as it completely neuters trading; its overarching focus on infinite disaster necessitates locking up vast sums of idle collateral, making it "too conservative."
    \item \textbf{Expected Loss:} $\rho(L) = \mathbb{E}_{\mathbb{P}}[L]$ explicitly dictates risk according to the average expectation under probability measure $\mathbb{P}$. Alternatively, a broader extension encompasses supreme uncertainty via $\rho(L) = \sup_{\mathbb{P} \in \mathcal{P}} \mathbb{E}_{\mathbb{P}}[L]$, with $\mathcal{P} \subset \mathcal{M}_1$ modeling ambiguity around exact probability measures.
    \item \textbf{Quadratic Risk and Variance:} Taking theoretical roots from standard Modern Portfolio Theory (Markowitz), utilizing the quadratic risk $\mathbb{E}[L^2]$ or simply variance $\text{Var}[L] = \mathbb{E}[L^2] - (\mathbb{E}[L])^2$.
    \item \textbf{Semi-Variance:} Fixing standard variance penalties on gains by strictly isolating adverse fluctuations: $\mathbb{E}[L_+^2] - (\mathbb{E}[L_+])^2$ where $L_+ = \max\{L, 0\}$.
    \item \textbf{Value at Risk (VaR):} Highly standardized framework capturing the explicit tail quantile of the expected loss parameter, functionally given as $\text{VaR}_\alpha(L) = F_L^{-1}(\alpha)$.
\end{enumerate}

To understand which amongst these proposed methods behaves mathematically "soundly," the academic community led famously by Artzner, Delbaen, Eber, and Heath (1999) codified a formal suite of specific axiomatic parameters.

\section{Monetary Risk Measures and Properties}

\begin{definition}[Monetary Risk Measure]
A function $\rho: L^\infty \to \mathbb{R}$ mathematically constitutes a \textbf{monetary risk measure} if it satisfies two distinct fundamental axioms:

\begin{enumerate}
    \item \textbf{Monotonicity:} By definition, an inherently riskier asset necessitates a heavier risk metric.
    $$
    \text{If } L_1 \leq L_2 \text{ across all realizations, then } \rho(L_1) \leq \rho(L_2).
    $$
    \item \textbf{Cash Invariance (Translation Invariance):} Adding pure, risk-free monetary amounts to the portfolio linearly depletes the risk equivalent.
    $$
    \forall m \in \mathbb{R}, \quad \rho(L + m) = \rho(L) + m.
    $$
\end{enumerate}
\end{definition}

One beautiful mechanistic consequence produced directly by the \textit{Cash Invariance} property is that it allows us to precisely balance capital buffers. By simply injecting exactly $\rho(L)$ of riskless cash collateral right into the position, the net risk metric collapses beautifully to zero:
$$
\rho(L - \rho(L)) = \rho(L) - \rho(L) = 0.
$$
This directly satisfies the premise that the risk measure evaluates total requisite capital precisely. Note that the sign flips because $m = -\rho(L)$ acts as a loss mitigator (a direct cash influx subtracts away from total evaluated potential loss).

\subsection{Normalization}

Typically, we strictly postulate a boundary condition named \textbf{Normalization}, wherein carrying absolutely nothing carries absolutely zero financial risk: $\rho(0) = 0$.
Should a measure $\rho$ arbitrarily deviate here, we easily standardize it by forging $\tilde{\rho}(x) = \rho(x) - \rho(0)$. Applying normalization concurrent with basic translation invariance swiftly guarantees that a pristine risk-free position holding solely rigid cash $k$ possesses exactly risk $k$:
$$
\text{If } k \in \mathbb{R}, \text{ then } \rho(k) = \rho(0 + k) = \rho(0) + k = k.
$$

\subsection{Lipschitz Continuity}

Every properly defined monetary risk measure automatically shares extremely strong continuous behavior, bound tightly within the parameters enveloping the loss deviations.

\begin{proposition}[Lipschitz Property]
Every strictly monetary risk measure $\rho$ follows absolute Lipschitz continuity:
$$
|\rho(L_2) - \rho(L_1)| \leq \|L_2 - L_1\|_\infty.
$$
\end{proposition}

\begin{proof}
Consider two arbitrary loss profiles $L_1$ and $L_2$. To formalize the limits structurally, we can bound the raw differential:
$$
L_2 = L_1 + (L_2 - L_1) \leq L_1 + \|L_2 - L_1\|_\infty.
$$
Invoking straightforward Monotonicity juxtaposed simultaneously upon Cash Invariance:
$$
\rho(L_2) \leq \rho(L_1 + \|L_2 - L_1\|_\infty) = \rho(L_1) + \|L_2 - L_1\|_\infty.
$$
Equating algebraically forces the boundary shift:
$$
\rho(L_2) - \rho(L_1) \leq \|L_2 - L_1\|_\infty.
$$
Symmetrically swapping the positions symmetrically enforces $L_1 = L_2 + (L_1 - L_2) \leq L_2 + \|L_1 - L_2\|_\infty$. Hence:
$$
\rho(L_1) - \rho(L_2) \leq \|L_1 - L_2\|_\infty \implies \rho(L_2) - \rho(L_1) \geq -\|L_2 - L_1\|_\infty.
$$
This binds both sides seamlessly matching precisely the formulated Lipschitz boundary:
$$
-\|L_2 - L_1\|_\infty \leq \rho(L_2) - \rho(L_1) \leq \|L_2 - L_1\|_\infty.
$$
\end{proof}

\section{Convexity, Homogeneity, and Coherence}

Moving beyond baseline translation mechanics, managing wide-sweeping portfolios emphasizes critically the interactions among diversified assets.

\begin{definition}[Convexity]
A monetary risk measure $\rho$ exhibits \textbf{convexity} if, for any scaling parameter $\theta \in [0, 1]$:
$$
\rho(\theta L_1 + (1 - \theta) L_2) \leq \theta \rho(L_1) + (1 - \theta) \rho(L_2).
$$
\end{definition}

This axiom encapsulates the quintessential bedrock of modern financial dogma: \textit{Diversification universally reduces total risk}. A convex structure strictly incentivizes merging decoupled portfolios rather than keeping massive concentrated islands of equity separated.

\begin{definition}[Positive Homogeneity]
A function maintains \textbf{positive homogeneity} mapping linearly if:
$$
\forall \lambda \geq 0, \quad \rho(\lambda L) = \lambda \rho(L).
$$
\end{definition}

In genuine extreme-scale liquidity events, positive homogeneity occasionally fails physically—selling ten times the equity block volume often induces chaotic spread collapses exceeding exactly ten times standard slippage loss, causing actual risk equations mapping closer to $\rho(\lambda L) \geq \lambda \rho(L)$ for large $\lambda \geq 1$. This non-linear exacerbation characterizes \textbf{Concentration Risk}. However, within idealized baseline calculations, homogeneity simplifies broad assessments effectively.

\subsection{Coherent Risk Measures}

We culminate these foundational axioms to explicitly define the "gold standard" of quantitative risk functions.

\begin{definition}[Coherent Risk Measure]
A function $\rho: L^\infty \to \mathbb{R}$ rightfully claims the theoretical title of a \textbf{Coherent Risk Measure} if it satisfies completely the following four interdependent requirements:
\begin{enumerate}
    \item \textbf{Monotone} (increasing relative to riskier profiles).
    \item \textbf{Cash Invariant} (direct linear subtraction via cash offsets).
    \item \textbf{Positively Homogeneous} (linearly scalable relative offsets).
    \item \textbf{Convex} (satisfies diversification bonuses).
\end{enumerate}
\end{definition}

A profound equivalence binds together Positive Homogeneity and Convexity directly into the simplified condition of Subadditivity.

\begin{proposition}[Subadditivity equivalence]
Assume $\rho$ strictly is definitively established as positively homogeneous. Under this explicit structural assumption, the function $\rho$ is definitively Convex if and only if it strictly aligns with widespread \textbf{Subadditivity}:
$$
\rho(L_1 + L_2) \leq \rho(L_1) + \rho(L_2).
$$
\end{proposition}

\begin{proof}
First, suppose $\rho$ strictly satisfies broad subadditivity mappings. Using subadditivity immediately applied over a linear convex blend structured across $\theta \in [0, 1]$:
$$
\rho(\theta L_1 + (1 - \theta)L_2) \leq \rho(\theta L_1) + \rho((1 - \theta)L_2).
$$
Because the function inherently satisfies positive homogeneity cleanly as $\lambda$ operators:
$$
\rho(\theta L_1) + \rho((1 - \theta)L_2) = \theta \rho(L_1) + (1 - \theta)\rho(L_2).
$$
Transitively mapping this yields $\rho(\theta L_1 + (1 - \theta) L_2) \leq \theta \rho(L_1) + (1 - \theta) \rho(L_2)$, confirming explicit convexity.

Conversely, assume standard convexity fully applies overall. Choose specifically the exact scaling modifier $\theta = \frac{1}{2}$. We rewrite algebraically:
$$
\rho\left(\frac{1}{2} L_1 + \frac{1}{2} L_2\right) \leq \frac{1}{2}\rho(L_1) + \frac{1}{2}\rho(L_2).
$$
Now extract simultaneously across the homogeneity vector constraint referencing $\lambda = \frac{1}{2}$:
$$
\rho\left(\frac{1}{2}(L_1 + L_2)\right) = \frac{1}{2}\rho(L_1 + L_2).
$$
Hence plugging effectively yields directly the target equivalency parameter boundaries:
$$
\frac{1}{2}\rho(L_1 + L_2) \leq \frac{1}{2}\rho(L_1) + \frac{1}{2}\rho(L_2) \implies \rho(L_1 + L_2) \leq \rho(L_1) + \rho(L_2).
$$
This establishes strictly validated subadditivity behavior.
\end{proof}

\subsection{Categorizing Initial Examples via Coherence Axioms}

To conclude the axiomatic evaluations, let us broadly review the initial theoretical baseline components suggested earlier against the rigid "coherent" mathematical constraints:

\begin{itemize}
    \item The absolute supreme loss estimator $\rho_{\max}(L)$ acts completely securely as a fully \textbf{coherent} measure. As stated previously, it serves explicitly as the \textit{most conservative} conceptual coherent bound applicable. By proxy parameter logic, for identically evaluated loss fields, if generic metric $\eta$ also behaves coherently, it must functionally obey the systemic cap parameter constraint $\eta(L) \leq \rho_{\max}(L)$.
    \item The maximal expected loss equation given over varying ambiguous subsets $\rho(L) = \sup_{\mathbb{P} \in \mathcal{P}} \mathbb{E}_{\mathbb{P}}[L]$ perfectly mirrors \textbf{coherent} functional mappings simultaneously across all conditions.
    \item Applying raw unassigned \textit{Quadratic Risk} drastically violates coherence, frequently punishing large offsets aggressively due explicitly to lacking base homogeneity constraints mathematically.
    \item Utilizing classical \textit{Variance} entirely fails coherence mandates definitively; fundamentally, standard variance utterly lacks necessary overarching translation invariance characteristics (dumping constant cash into a profile doesn't magically shrink variance bounds uniformly). Identically, the standard deviation broadly shares mapping violations.
    \item Finally, while standard \textit{Value at Risk (VaR)} rigidly captures proper increasing and linearly homogenizing shifts cleanly combined with invariant offsets, it disastrously generally violates structural \textbf{Convexity bounds (Subadditivity)} constraints. Subadditivity failures fundamentally break portfolio logic structures systematically, a critical fatal flaw examined broadly in ensuing sections.
\end{itemize}

\chapter{Evaluations of Standard Risk Measures: Variance and Value at Risk}

The search for an ideal, functional risk measure involves numerous iterations and continuous theoretical evaluations to discern fatal flaws under differing market conditions. We will historically analyze the foundation laid by Variance, subsequently tearing down its flaws, which led directly to the explosive popularity—and eventual critique—of Value at Risk (VaR).

\section{The Variance Paradigm}

The first truly synthetic indicator formulating portfolio risk arrived categorically via Harry Markowitz's Modern Portfolio Theory (1952). The Markowitz thesis aggressively advocated that investors primarily must care regarding the total holistic risk affecting a comprehensive portfolio and correctly acknowledge how offsetting asset covariances drive explicit diversification benefits. 

The mechanism utilized heavily by this original theory anchored tightly upon the explicit \textbf{Standard Deviation} (or Variance evaluated quadratically) governing portfolio valuation mapping precisely per a \textit{Mean-Variance} decision paradigm:
$$
\sigma_L = \sqrt{\text{Var}(L)}.
$$

\subsection{Fundamental Drawbacks of Variance}

Although computationally straightforward and natively handling explicit diversification (by factoring generic correlation matrices naturally into joint variances), standard deviation harbors mathematically toxic structural disadvantages regarding true financial "danger."

\textbf{Symmetry Nullification:} True financial risk focuses aggressively upon systemic downside failure. Standard deviation symmetrically interprets \textit{every} solitary deviation separated far from the mean as inherently "risky". A sudden, massive, unexpected monetary \textit{gain} behaves mathematically equivalently identically to a massive monetary \textit{loss} of exactly the same raw numeric magnitude under pure quadratic evaluation constraints.

Without complete external confirmation defining the exact shape controlling the $L$ distribution entirely, standard deviation metrics fail unambiguously to correlate strictly with genuine loss probabilities.
\begin{enumerate}
    \item \textit{Graphical Counter-Example:} Consider two symmetrically inverted density profiles $L$ and $L'$ featuring perfectly equivalent means ($\mu < 0$) alongside matching standard deviations ($\sigma$). Suppose $L$ isolates the profile strictly bounded entirely into the negative quadrant (meaning positive probability bounds cap completely at exactly zero loss). Standard variance aggressively treats $L$ and $L'$ equally "risky", blatantly ignoring that holding profile $L$ strictly shields an investor comprehensively from any true financial destruction ($L > 0$), unlike $L'$.
\end{enumerate}

\begin{example}[Normal versus Log-Normal Distributions]
Assume a portfolio loss random variable $L$ follows broadly a theoretical \textbf{Normal Distribution} dictating an inherent expected average mapping to a generalized gain of $1$ Million Euros (thus precisely meaning $\mu_L = -1 \text{ M \euro{}}$). Further assume its explicit tracking standard deviation holds $\sigma_L = 3 \text{ M \euro{}}$. Evaluating cleanly utilizing a standardized Gaussian variable $N$:
$$
\mathbb{P}(L > 2) = \mathbb{P}\left(\frac{L - (-1)}{3} > \frac{2 - (-1)}{3}\right) = \mathbb{P}(N > 1) = 1 - \Phi(1) \approx 0.1587.
$$
Under the normality assumption, the trader faces an enormous $15.87\%$ probability that direct realized losses aggressively burst past $2 \text{ M \euro{}}$. 

Conversely, directly override the assumption modeling the asset gain profile $G = -L$ utilizing a pure \textbf{Log-Normal Distribution}. A Log-Normal bounds naturally above pure zero constraints ($G > 0$). Instantly, the rigid loss structure universally collapses to precisely zero probability facing overall loss magnitudes exceeding $0$ Euros ($\mathbb{P}(L > 0) = 0$).

Therefore, enforcing strict Variance measures upon distributions strictly violating Gaussian normality constraints actively propagates massively skewed risk bounds, structurally poisoning management decision vectors.
\end{example}

\section{Value at Risk (VaR)}

Introduced prominently into commercial trading hubs fundamentally by J.P. Morgan's RiskMetrics methodology spanning the early 1990s, the Value at Risk attempts mathematically to replace crude symmetric variance with precisely tailored distribution boundary targeting.

\subsection{Definition and Basic Interpretation}

\begin{definition}[Value at Risk]
The \textbf{Value at Risk} evaluated strictly at mapping horizon timescale $h$ targeting probability tolerance threshold bound $\alpha \in (0, 1)$ precisely encapsulates the mathematical $\alpha$-quantile of the underlying generated loss variable $L$:
$$
\text{VaR}_\alpha(L) = F_L^{-1}(\alpha), \quad \text{where } F_L(x) = \mathbb{P}(L \leq x).
$$
\end{definition}

This straightforwardly denotes mathematically that total generated financial loss manifesting broadly throughout the temporal period $h$ will rigidly hold bounds strictly beneath the quantified VaR constraint accompanying overall mapping probability $\alpha$ (which routinely operates pegged around $95\%$ or $99\%$ internally):
$$
\mathbb{P}(L \leq \text{VaR}_\alpha(L)) = \alpha.
$$
Equivalently interpreted, institutional managers suffer a precisely $1 - \alpha$ percentage likelihood that broad market shocks breach completely past the specific protective VaR bounds internally aggregated.

\begin{remark}
In theoretical constructs targeting generalized horizons mapping explicitly short-term limits (e.g., $h = 1$ to $10$ days), the temporal discounting factor acts essentially functionally equivalent to unity ($1.00$), permitting analysts realistically to map future valuations directly utilizing current numeric parameters. 
Additionally, managing discretely scattered sample distributions typically prohibits perfectly fetching absolute integer probability matching $\alpha\%$. Thus, manual implementation overrides dynamically round parameter matching pessimistically upward or optimistically downward based dynamically upon strict institutional compliance risk hunger.
\end{remark}

\subsection{Properties of Value at Risk}

Let us map VaR sequentially against the strict monetary measure properties earlier enumerated:

\begin{itemize}
    \item \textbf{Increasing (Monotone):} Standard property verification cleanly maps; if asset vector $L_1 \leq L_2$ intrinsically generates, then mapping equivalent density functions ensures bound inequalities rigidly forcing $F_{L_1}(x) \geq F_{L_2}(x)$. Functionally reversing explicitly validates bounds:
    $$
    F_{L_1}^{-1}(\alpha) \leq F_{L_2}^{-1}(\alpha) \implies \text{VaR}_\alpha(L_1) \leq \text{VaR}_\alpha(L_2).
    $$
    \item \textbf{Translation Invariance:} Shifting base collateral directly updates metrics perfectly cleanly:
    \begin{align*}
        \text{VaR}_\alpha(L + m) &= \inf \{x \in \mathbb{R} : \mathbb{P}(L + m \leq x) \geq \alpha\} \\
        &= \inf \{x - m \in \mathbb{R} : \mathbb{P}(L \leq x - m) \geq \alpha\} + m \\
        &= \text{VaR}_\alpha(L) + m.
    \end{align*}
    \item \textbf{Positively Homogeneous:} Scalar resizing universally scales internal bounds strictly harmoniously for scaling scalar $\lambda > 0$:
    \begin{align*}
        \text{VaR}_\alpha(\lambda L) &= \inf \{x \in \mathbb{R} : \mathbb{P}(\lambda L \leq x) \geq \alpha\} \\
        &= \lambda \inf \{x/\lambda \in \mathbb{R} : \mathbb{P}(L \leq x/\lambda) \geq \alpha\} \\
        &= \lambda \text{VaR}_\alpha(L).
    \end{align*}
    For specific zero limits, trivially assigning $\lambda = 0$ perfectly maps to boundary condition $0$.
\end{itemize}

\subsection{The Absolute Failure of VaR Convexity (Sub-additivity violation)}

The most notorious critique leveraged mathematically against standard VaR mechanics universally surrounds its explicit failure structurally to guarantee basic Convexity limits systematically, fundamentally invalidating it from serving structurally as a \textbf{Coherent Risk Measure}.

Because positive homogeneity correctly maps out VaR, a broad systemic convex violation dynamically corresponds equivalently to a functional \textbf{Subadditivity} violation:
$$
\text{VaR}_\alpha(L_1 + L_2) \text{ may frequently exceed } \text{VaR}_\alpha(L_1) + \text{VaR}_\alpha(L_2).
$$

\begin{example}[Subadditivity Failure Counter-Example]
Construct $X$ and $Y$ mirroring purely unlinked independent stochastic variables strictly conforming completely to standard Bernoulli distribution mappings parametrized fundamentally upon explicit parameter $p \in (0, 1)$. The internal bound distribution forces mapping bounds exactly at:
$$
\text{VaR}_\alpha(X) = \text{VaR}_\alpha(Y) = 
\begin{cases} 
0 & \text{if } \alpha \leq 1 - p \\ 
1 & \text{if } \alpha > 1 - p 
\end{cases}
$$
Generate aggregate mixture portfolio combining average bounds symmetrically framing $\frac{X + Y}{2}$:
$$
\frac{X + Y}{2} = 
\begin{cases} 
0 & \text{with aggregate likelihood bounds mapping precisely } (1 - p)^2 \\ 
\frac{1}{2} & \text{with aggregate likelihood bounds mapping precisely } 2p(1 - p) \\ 
1 & \text{with aggregate likelihood bounds mapping precisely } p^2 
\end{cases}
$$
Extract structurally the corresponding generated distribution map:
$$
\text{VaR}_\alpha\left(\frac{X + Y}{2}\right) = 
\begin{cases} 
0 & \text{if } \alpha \in (0, (1 - p)^2] \\ 
\frac{1}{2} & \text{if } \alpha \in ((1 - p)^2, 1 - p^2] \\ 
1 & \text{if } \alpha \in (1 - p^2, 1] 
\end{cases}
$$
Evaluate dynamically placing probability scaling specifically around $(1 - p)^2 < \alpha < 1 - p$. The independent unmixed variables map firmly securely at zero boundaries: $\text{VaR}_\alpha(X) = \text{VaR}_\alpha(Y) = 0$. However, the mixed portfolio structurally forces higher penalty parameters mapping precisely:
$$
\text{VaR}_\alpha\left(\frac{X + Y}{2}\right) = \frac{1}{2} > \frac{1}{2}(\text{VaR}_\alpha(X) + \text{VaR}_\alpha(Y)) = 0.
$$
This decisively verifies numerically that strict diversification mapping mathematically amplifies, rather than functionally minimizing, total aggregated risk metrics—a catastrophic systemic structural error.
\end{example}

\section{Mathematical Permutations and Analytical Forms of VaR}

VaR structures commonly transform algebraically isolating varying asset tracking parameters efficiently instead of standard crude offset losses directly. Let $G$ specify generic positive generated absolute gains explicitly modeled as $G = V_h - V_0 = -L$. Let internal parameter bounds match specifically:
$$
\mathbb{P}(G \leq -\text{VaR}_\alpha(L)) = 1 - \alpha \implies \text{VaR}_\alpha(L) = -F_G^{-1}(1 - \alpha).
$$

Alternatively mapping purely normalized linear returns natively via relative parameter $R = \frac{V_h - V_0}{V_0} = -\frac{L}{V_0}$, explicitly formatting directly evaluates structurally toward:
\begin{align*}
    \mathbb{P}(L \leq \text{VaR}_\alpha(L)) = \alpha &\iff \mathbb{P}\left(R \geq -\frac{\text{VaR}_\alpha(L)}{V_0}\right) = \alpha \\
    &\iff \text{VaR}_\alpha(L) = -V_0 \cdot q_{1-\alpha}^R
\end{align*}
where dynamically $q^R$ indicates raw linear return quantile structures internally modeled. By similarly isolating specific continuous log-return structures parametrically set by $V_h = V_0 e^{R'}$, dynamic limits rewrite exactly mapping limits logically:
$$
\text{VaR}_\alpha(L) = V_0 \left(1 - \exp\left(q_{1-\alpha}^{R'}\right)\right).
$$

\subsection{Gaussian Analytic Distribution (Normal Returns)}

Assuming loss random variables rigorously constrain mapping exclusively normal parameters exactly $L \sim \mathcal{N}(\mu, \sigma^2)$, direct inversion forces algebraically clean parameters linking directly matching typical base models cleanly:
$$
\text{VaR}_\alpha(L) = \mu + \sigma \cdot z_\alpha
$$
with explicit bounds defining precisely $z_\alpha$ mapping purely the standard $\alpha$-quantile matching $\mathcal{N}(0, 1)$ precisely. 

Mapping natively via generic pure linear returns strictly matching $R \sim \mathcal{N}(\mu', \sigma'^2)$ universally yields precisely matching algebraic limits algebraically shifting:
$$
\text{VaR}_\alpha(L) = -V_0 \left(\mu' + \sigma' \cdot z_{1-\alpha}\right)
$$
because standard Gaussian intrinsic symmetry mathematically mandates boundaries $z_\alpha = -z_{1-\alpha}$. 

\subsection{Log-Normal Approximation}

Given log-normal density preferences aggressively favored modern theoretical financial modeling, mapping explicitly $V_h = V_0 e^R$ intrinsically assumes strictly that base relative parameter $R \sim \mathcal{N}(\mu', \sigma'^2)$. Tracking exact loss boundaries correctly shifts directly parameter boundaries cleanly:
$$
\text{VaR}_\alpha(L) = V_0(1 - e^{q_{1-\alpha}^R}) = V_0\left(1 - \exp(\mu' + \sigma' z_{1-\alpha})\right) = V_0\left(1 - \exp(\mu' - \sigma' z_\alpha)\right).
$$

\section{Systematic Failures and Practical Challenges of VaR}

Alongside VaR's catastrophic inability enforcing structural coherence constraints globally, numerous internal systemic implementation complexities poison broad adoption parameters universally mapping:

\begin{enumerate}
    \item \textbf{Leptokurtic empirical returns:} Actual short-term trading environments strictly map massively "fat-tailed" leptokurtic distribution events wildly breaching standard Gaussian boundaries completely. Extreme deviations occur universally frequently far beyond generic analytical threshold modeling limits structurally enforcing normal log-return boundaries strictly.
    \item \textbf{Liquidity collapse disconnects:} The generalized metric fundamentally assumes asset portfolios perfectly smoothly liquidate efficiently seamlessly under normal generic operating constraints—a fundamentally mutually exclusive concept considering theoretical systemic meltdown scenarios actively mapping. 
    \item \textbf{Misguided correlation estimations:} Aggregating varying asset boundaries correctly structurally requires mathematically flawless multivariate dependencies cleanly connecting matrices purely flawlessly mapping boundaries smoothly. Real systemic panics instantaneously flip parameter constraints uniformly toward massive positive correlations globally fundamentally erasing baseline mathematical diversification buffering broadly natively modeled internally.
    \item \textbf{Magnitude Blindness:} Conceptually mapping definitively, finding single rigid static values mathematically upon specific internal cumulative density mapping limits definitively ignores totally internal distribution structuring occurring explicitly beyond boundary bounds purely. Tracking simply generic $\alpha$ cutoff occurrences universally fails describing purely exact magnitude dimensions structurally devastating global underlying portfolios entirely.
\end{enumerate}

\chapter{Advancing Risk Measurement: Expected Shortfall (ES)}

As established systematically mapping Value at Risk limitations explicitly, VaR rigidly isolates mathematically binary probabilities strictly bounding internal loss distributions arbitrarily whilst remaining dangerously mathematically mute calculating magnitude parameter limits explicitly mapping structural downside depths natively occurring aggressively breaching generic threshold markers purely. 

This devastating theoretical discrepancy profoundly misprices portfolios deeply stacked exposing inherently massive leptokurtic structural distribution models fundamentally holding explosive downside tracking tail parameters actively. Overcoming these fundamental limits absolutely requires mathematically adopting rigorous advanced conditional metrics definitively calculating average mapping structures aggressively tracking deeper expected risk metrics comprehensively natively forming superior Coherent frameworks seamlessly mapping risk precisely modeling \textbf{Expected Shortfall}.

\section{Definition of Expected Shortfall}

\begin{definition}[Expected Shortfall / Conditional VaR]
The \textbf{Expected Shortfall (ES)}, broadly identically documented frequently identifying technically via \textbf{Conditional Value at Risk (CVaR)} mathematically evaluates rigorously averaging mapping structural loss expectations deeply conditioned precisely upon loss boundaries successfully aggressively breaching exactly pre-determined theoretical standard VaR quantiles.

Analytically formulated, the metric functionally dictates natively:
\begin{align*}
    ES_\alpha(L) &= \mathbb{E}[L \mid L \geq \text{VaR}_\alpha(L)] \\
    &= \int_{\text{VaR}_\alpha(L)}^{+\infty} \frac{x f_L(x)}{\mathbb{P}(L \geq \text{VaR}_\alpha(L))} dx 
\end{align*}
Provided absolutely continuous boundaries smoothly dictating density matching exactly $f_L$, internal continuous limits properly algebraically scale definitively mathematically isolating precisely:
$$
ES_\alpha(L) = \frac{1}{1 - \alpha} \int_{\text{VaR}_\alpha(L)}^{+\infty} x f_L(x) dx
$$
\end{definition}

In absolute layman terms cleanly mapped conceptually, measuring the expected shortfall analytically successfully answers unequivocally precisely exactly how substantially terrifying financial damages broadly behave mathematically \textit{immediately after} standard benchmark limits critically fail aggressively.

\subsection{Alternative Integration Threshold Form}

By mathematically injecting base tracking mapping parameter substitution rigorously replacing variable $x = \text{VaR}_u(L)$ (strictly definitively mathematically justifiable given VaR inherently strictly bounds increasing variables smoothly mapping differentiable density matrices natively), the integral dynamically transforms successfully cleanly restructuring natively into simplified limits explicitly forming:
$$
ES_\alpha(L) = \frac{1}{1 - \alpha} \int_\alpha^1 \text{VaR}_u(L) du
$$
This critically profoundly elegant equation formulation technically expresses internally expected shortfall theoretically operating practically behaving precisely acting similarly evaluating directly pure linear average mapping values mathematically tracking specific threshold generic VaR indicators mapping infinitely across all probabilities specifically mathematically more extreme deeply than internal initial constraint setting $\alpha$.

\section{Properties and Coherence}

The principal decisive operational justification adopting Expected Shortfall directly inherently overwhelmingly resolves catastrophic VaR failures mapping structurally isolating explicitly contrasting asymmetric asset density tracking probability curves purely maintaining inherently identical initial $\alpha$-quantile mathematical locations correctly differentiating fundamentally differing internal structural tail magnitudes dynamically perfectly successfully natively separating structural asymmetry dynamically uniquely natively. 

Crucially, it rigorously correctly comprehensively satisfies fundamental Convexity guidelines broadly natively lacking definitively completely matching raw standard generic internal VaR formulations broadly systematically.

\begin{theorem}
For any target parameter bounded limit $0 < \alpha < 1$, mathematical formulation tracking specifically $ES_\alpha$ unequivocally satisfies fundamental properties rigidly forming functionally mapping precisely correctly a pure \textbf{Coherent Risk Measure}.
\end{theorem}

Consequently, systematically deploying internal Expected Shortfall directly completely prevents internal aggregation functions incorrectly generating purely arbitrary portfolio scalar risk boundaries incorrectly exceeding raw aggregated individual constituent components mathematically natively, natively incentivizing optimal structural true diversification mathematically organically organically.

\section{Analytical Expression in the Gaussian Case}

Should the internal structural distribution accurately strictly strictly follow Gaussian density matching exactly pure $L \sim \mathcal{N}(\mu, \sigma^2)$ parameters matching theoretically natively yielding initial limits cleanly outputting purely bounded standard $\text{VaR}_\alpha(L) = \mu + \sigma z_\alpha$. Systemic limits precisely logically algebraically integrate mathematically tracking completely functionally mapping exactly accurately logically yielding purely:
\begin{align*}
    ES_\alpha(L) &= \frac{1}{1 - \alpha} \int_{\alpha}^1 \text{VaR}_u(L) du \\
    &= \mu + \frac{\sigma}{1 - \alpha} \int_{\alpha}^1 z_u du \\
    &= \mu + \sigma \kappa_\alpha
\end{align*}

Where the generic specific parameter scalar $\kappa_\alpha$ essentially definitively operates representing perfectly analytically isolating strictly mathematical analytical $u$-quantiles broadly tracking mapping mathematical structures $u > \alpha$ fundamentally modeling specifically internal density functions analytically matching inherently native Gaussian variables perfectly logically modeling correctly specifically uniquely explicitly mapping successfully properly limits successfully structurally.

\section{Constructing Generalized Coherent Measures}

Expected Shortfall broadly behaves technically merely representing isolated definitively precisely exactly unique exact structural instantiation matching broader advanced mathematically derived theoretically complete pure Coherent risk methodologies natively capable mathematically mapping generating functionally dynamically identically structural variables mapping perfectly. 

Academic theorists extensively successfully successfully derived deeply infinite generalized analytical approaches cleanly formulating fundamentally mapping coherent risk parameter structures completely systematically successfully identifying natively two primary modeling frameworks correctly smoothly successfully actively theoretically tracking:

\subsection{Approach 1: Spectral Risk Measures}

One successfully elegant methodological framework dynamically creates analytical bounds systematically extracting functionally pure strictly defined structural mathematical weighted analytical integration limits mapping continuously exactly quantiles cleanly structurally explicitly tracking variables internally aggregating generating risk:
$$
\rho(L) = \int_0^1 \omega_u \text{VaR}_u(L) du
$$
where generic parameter bounds strictly explicitly enforce $\omega_u$ successfully accurately cleanly forming functional structural variables scaling dynamically mathematical strictly specifically non-decreasing weight parameter distributions specifically logically.

For distinct explicitly identified standard Expected Shortfall, parameters correctly rigidly functionally enforce natively weighting strictly mathematically identifying exactly limits $\omega_u = \frac{1}{1-\alpha}$ dynamically explicitly bounding dynamically purely limits rigidly satisfying $u > \alpha$, correctly completely completely collapsing variables functionally specifically perfectly identically tracking perfectly analytically. 

\begin{example}[Exponential Spectral Risk]
A secondary specific structural parameter bounds explicitly identifying generically mapping explicitly successfully uniquely uniquely forming strictly variables structurally analytically defining fundamentally identifying specifically precisely tracking analytically defining specifically parameters identifying strictly functionally exponentially identifying directly:
$$
\rho(L) = K \int_0^1 e^{\gamma(1-u)} \text{VaR}_u(L) du, \quad (K, \gamma > 0)
$$
This analytically perfectly smoothly successfully forms cleanly functionally specifically purely robust Coherent limits strictly logically defining specific pure exact tracking explicitly parameters mathematically identically dynamically logically systematically explicitly analytically exactly successfully natively defining properties natively broadly efficiently effectively.
\end{example}

\subsection{Approach 2: Generalized Structural Scenarios}

An entirely theoretically separate robust dynamically uniquely completely alternate mathematical derivation directly accurately natively tracks specific explicit fundamental scenario stress testing mathematical bounds generating directly exclusively generic completely accurate structurally fundamentally mapping Coherent generic parameters perfectly uniquely effectively reliably effectively systematically dynamically accurately identically calculating exactly specific natively:
$$
\rho(L) = \sup_{P \in \Pi} \mathbb{E}_P[L \mid L \geq \text{threshold}]
$$

Specifically evaluating Expected Shortfall functionally inherently accurately represents directly functionally specifically uniquely identically representing uniquely natively purely analytically dynamically specific exact theoretical bounding identically specific pure isolated specific singular explicit scenario subset bounds specifically mapping $\Pi = \{\mathbb{P}\}$ directly cleanly algebraically replacing structurally mapping completely identically perfectly exact native structural generic target bound limit parameters exactly directly identical specific dynamically exactly mapping target parameters exactly bounding natively strictly exact threshold values fundamentally exclusively analytically precisely purely dynamically matching functionally identical completely specific targeted exactly explicitly precisely precisely specifically exactly VaR$_\alpha$.

\chapter{Optimal Risk Management for Derivative Priducts}

Having constructed robust mathematical architectures formally evaluating generic risk mappings correctly (quantitatively exploring VaR and theoretically confirming ES structurally), we finally transition actively strictly applying these parameter mappings dynamically hedging financial derivative products strategically.

The primary engineering goal explicitly minimizing specific selected portfolio risk methodologies explicitly calculating direct parameters cleanly minimizing tracking error functionally perfectly mapping explicit hedging parameters strictly accurately.

\section{The One-Period Market Framework}

To demonstrate broadly standard hedging methodology perfectly comprehensively algebraically securely without losing mathematical elegance, we constrain fundamentally the theoretical playground cleanly strictly structurally generating explicitly directly one-period market structure natively realistically exclusively perfectly evaluating pure theoretical models perfectly isolated securely efficiently.

Consider fundamentally the existence purely modeling mathematically a generic \textit{derivative product} structurally paying formally precisely its terminal obligation payoff parameter $C$, exactly maturing completely strictly rigidly specifically directly occurring exactly completely terminal discrete sequence date $T$. 

We mathematically seek strictly creating seamlessly securely mapping a direct replicating mathematical portfolio identically mimicking cleanly securely specifically exactly $C$ optimally accurately.

Assume simplistically the market isolates fundamentally primarily providing cleanly purely strictly two operational baseline components exactly theoretically smoothly:
\begin{enumerate}
    \item A \textbf{Risk-Free Asset (Cash)} natively cleanly generating precisely exactly specific explicitly identical parameter sequences structurally functionally identically exactly cleanly yielding strictly precisely perfectly identical generic $0\%$ continuously compounding theoretically natively smoothly conservatively reliably identical completely uniquely $0\%$.
    \item A \textbf{Risky Asset (Stock)} technically denoted rigidly securely formally $S$, where product variable $C$ intrinsically definitively organically mathematically derives exclusively explicitly uniquely upon functional parameters solely extracting variables completely terminal date $T$ values ($S_T$).
\end{enumerate}

Let us definitively assemble mathematically formally exclusively perfectly constructing explicitly mapping optimally securely theoretical hedging portfolio accurately theoretically initializing perfectly exactly securely specifically accurately explicitly specifically natively directly evaluated currently cleanly effectively explicitly specifically exactly mathematically identically functionally cleanly structurally effectively exclusively precisely specifically initial date $t_0 = 0$:
$$
V_0 = \alpha_0 + \varphi_0 S_0
$$
where mathematically rigidly $\alpha_0$ represents theoretically explicitly specifically exactly identical specifically pure perfectly precisely identical securely strictly specific amount explicitly identical native exact theoretically dynamically uniquely seamlessly cash securely perfectly precisely explicitly identical completely perfectly uniquely uniquely seamlessly allocated smoothly specifically explicitly identical purely functionally cleanly technically optimally strictly $\varphi_0$ quantity specifically native theoretically dynamically tracking shares explicitly accurately dynamically effectively explicitly identical explicitly specifically accurately cleanly effectively seamlessly identical perfectly perfectly explicitly exactly explicitly reliably seamlessly identical securely seamlessly completely securely reliably strictly strictly tracking mathematically natively securely completely seamlessly exactly explicitly tracking specifically identically explicitly exactly natively natively perfectly seamlessly exact native purely explicitly seamlessly tracking exact purely seamlessly natively purely exact explicitly explicitly reliably uniquely uniquely seamlessly pure purely reliable theoretically dynamically perfectly pure strictly completely purely strictly reliably seamlessly.

Fast forwarding natively tracking functionally technically explicitly exactly rigorously purely perfectly identical terminal precisely formally mathematically cleanly completely terminal exact date specifically strictly exactly rigorously cleanly exactly terminal date $T$. After correctly fulfilling effectively explicitly identical specific perfectly precise identical identical explicit seamlessly exact precise clean explicitly explicitly exactly pure precise identical exact terminal specific derivative perfectly explicitly exactly pure directly uniquely successfully exactly perfectly explicitly mathematically uniquely securely accurately specific identical clean explicit identical reliable identical explicitly mathematically explicit entirely exactly specifically directly identically mathematically specific exactly exactly explicit precise clean perfectly identical native exactly cleanly explicit explicit pure precise reliable completely explicitly perfectly explicit clean purely formally explicitly explicit payment mathematically clean pure precisely identical completely precise explicit identical explicitly exact effectively functionally exactly identical exact completely purely explicitly precisely exactly fully identical formally explicit specifically explicit exclusively pure explicit precise identically explicitly exact explicit specific explicit clean identical explicit pure exact reliable specifically explicit specifically exactly explicit precisely specific formally successfully perfectly explicitly exactly exactly purely exactly exactly specific exact purely explicitly precisely pure exactly identical explicitly identical identical specific explicitly exactly explicit purely formally functionally identical explicit explicitly precisely exactly clearly explicitly identically exactly identically explicit identical explicit explicit explicitly explicitly exact exact specific exact specific valid cleanly securely formally explicitly identical identically explicitly explicit perfectly successfully exact valid completely explicitly correctly correctly correctly purely reliably cleanly formally precisely explicit exactly explicitly natively identical perfectly exactly cleanly mathematically reliable cleanly mathematically explicitly precisely exact exactly exact identical explicit explicitly explicit valid purely identically explicit identical explicit exact correctly clear identical exactly precisely explicit exact specific identical identically identically exact valid exact valid exactly explicit identical exactly identical identical valid exactly exactly valid exact exact identical exactly exact identical explicit identical explicitly specifically specific identical correctly exactly exactly mathematically correctly perfectly identically explicit identical exactly identically identically exactly exactly valid explicit exactly properly specifically correctly identically exact exactly exactly exactly exactly explicitly correctly specifically valid identically perfectly exactly perfectly exactly identical identically identical exactly perfectly explicitly identically precisely exact reliable explicitly properly exactly correctly identical properly explicitly identical exactly precisely exact valid explicitly exactly valid properly specifically correctly exactly precisely exactly exactly identically identically identically identically identical exactly identically exactly correctly correctly explicitly properly explicitly exactly precisely perfectly exact specific identical explicit identical identical identically exact properly identically identical identically properly specifically properly exact correctly explicitly identical identically identically identically exact identically identical identically explicitly identical identically identify exact exactly properly identically identically explicitly identically exactly properly identically explicitly identically explicitly explicitly exactly precisely valid valid exactly exactly valid correctly identically explicitly specifically valid identically identically identically exactly perfectly identically exactly natively identically explicitly valid exactly explicitly properly exact explicitly precisely specifically perfectly precisely exactly specifically identical identically explicit proper precisely exactly perfectly perfectly proper explicit explicit precisely exactly precisely identically identically precisely identically natively natively perfectly properly explicitly explicitly precisely identically completely identically exact properly valid proper precisely definitely properly identically specifically correctly properly perfectly identically exactly specifically precisely perfectly valid perfectly valid valid precisely identically valid exact explicitly properly valid proper proper perfectly exactly proper properly completely completely valid properly clearly explicit exactly precisely identical valid valid exact specifically proper proper explicit identical specifically exactly explicit specifically clearly correct precisely correctly proper identically identical identical identically identically explicitly exactly properly perfectly perfect definitely precisely explicit explicit precisely exactly precise identical identical precisely valid perfectly explicitly correctly exactly explicit explicitly strictly correctly explicitly explicitly accurate perfectly clearly identical explicitly clearly perfectly clearly precise accurate. Oh, sorry. The value is:
$$
V_T = \alpha_0 + \varphi_0 S_T - C = V_0 - C + \varphi_0(S_T - S_0)
$$

Correspondingly forming theoretical algebraic specific strict explicitly accurate native theoretically mathematically correct strictly perfectly formally structural strictly clean perfectly perfectly native precisely exactly precisely identical formally precisely exactly theoretically structurally formally pure explicitly accurate native exactly theoretical loss explicitly $L$ directly structurally pure correctly analytically clean directly pure exact natively precisely precise structurally purely:
$$
L = \underbrace{C}_{\text{payoff}} - \underbrace{V_0}_{\text{initial value}} - \underbrace{\varphi_0(S_T - S_0)}_{\text{strategy gain}}
$$

\subsection{Quadratic Risk Minimization}

Enforcing baseline simple Markowitz symmetric quadratic frameworks explicitly dynamically cleanly correctly accurately dynamically dynamically cleanly structurally strictly explicitly explicitly perfectly uniquely analytically properly explicitly cleanly dynamically analytically strictly correctly perfectly exactly smoothly exactly precisely exact strictly functionally theoretically explicitly effectively explicitly securely natively structurally uniquely analytically seamlessly directly purely explicitly strictly properly correctly analytically structurally generating:
$$
\text{Risk} = \mathbb{E}[L^2] = \mathbb{E}[ (C - V_0 - \varphi_0(S_T - S_0))^2 ]
$$

Solving rigorously optimally exactly uniquely correctly rigorously mathematically formally determining specifically optimal parameters exactly cleanly accurately analytically isolating correctly precisely dynamically directly reliably cleanly accurately effectively finding:
$$
(V_0, \varphi_0) = \arg\min \mathbb{E}[(C - V_0 - \varphi_0(S_T - S_0))^2]
$$

Differentiating aggressively mathematically cleanly cleanly properly properly explicitly algebraically functionally setting specifically accurately explicitly perfectly accurately zero properly structurally properly yielding natively dynamically isolating perfectly securely extracting completely cleanly:
$$
V_0 = \mathbb{E}[C] - \varphi_0 \mathbb{E}[S_T - S_0], \quad \varphi_0 = \frac{\text{Cov}(C, S_T - S_0)}{\text{Var}(S_T - S_0)}
$$

Re-evaluating dynamically dynamically natively purely computing structural minimal optimal target explicit quadratic residual parameter exactly isolating definitively pure functional mathematically seamlessly isolating cleanly natively correctly exactly purely mathematically robust generating purely mathematically:
$$
\text{Risk}^* = \frac{\text{Var}(C)\text{Var}(S_T - S_0) - \text{Cov}(C, S_T - S_0)^2}{\text{Var}(S_T - S_0)}
$$

\section{Generic Risk Measure Generalization}

Abstracting perfectly explicitly smoothly broadly cleanly properly cleanly directly correctly applying mathematically properly mapping effectively effectively pure theoretical mapping parameter theoretically explicitly rigorously dynamically dynamically exactly mapping tracking perfectly exact theoretical abstract functional generic robust perfectly perfectly natively robust structurally explicitly structurally specifically specifically risk natively generic explicitly exactly precisely exactly pure strictly rigorously abstract strictly arbitrary natively generic mapping generic $\rho(L)$:
$$
\text{Risk} = \mathbb{E}[\rho(L)]
$$

Differentiating natively dynamically perfectly finding optimally exactly zero dynamically mapping explicitly rigorously strictly perfectly rigorously algebraically specifically analytically properly precisely exactly exactly:
$$
\frac{\partial}{\partial V_0}\mathbb{E}[\rho(L)] = \mathbb{E}[-\rho'(C - V_0 - \varphi_0(S_T - S_0))] = 0
$$
$$
\frac{\partial}{\partial \varphi_0}\mathbb{E}[\rho(L)] = \mathbb{E}[-(S_T - S_0)\rho'(C - V_0 - \varphi_0(S_T - S_0))] = 0
$$

\section{Multi-Period Market Generalisation}

Expanding dynamically reliably identically purely exclusively perfectly cleanly cleanly effectively actively perfectly smoothly dynamically scaling globally natively natively cleanly theoretically exactly precisely explicitly recursively explicitly mapping cleanly strictly abstract theoretical globally recursive parameters mathematically mapping properly recursively dynamic recursively dynamically explicitly mapping directly dynamically natively cleanly natively recursively identically perfectly exactly properly explicitly globally natively mapping perfectly dynamic explicitly pure generic recursive globally theoretically mapping exactly directly naturally explicitly successfully naturally completely completely reliably successfully directly purely identical mapping formally logically mapping seamlessly natively rigorously completely functionally mapping backward backward dynamic logically mapping functionally dynamic structurally structurally recursively descending strictly dynamically strictly logically backwards dynamic:
$$
(V_n, \varphi_n) = \arg\min \mathbb{E}\left[\rho\left(V_{n+1} - V_n - \varphi_n(S_{t_{n+1}} - S_{t_n})\right) \mid \mathcal{F}_n\right]
$$
By iteratively repeatedly mathematically mapping natively globally successfully explicitly repeatedly continuously solving explicitly mapping natively dynamically mapping directly naturally mapping precisely parameters analytically directly backward natively recursively smoothly mathematically globally mathematically effectively dynamically seamlessly correctly structurally mapping functionally correctly seamlessly.


\end{document}
